\section{今後の課題}

本研究では、エッジ環境（Webブラウザ）上でリアルタイムかつ高精度に待機列およびそのグループ構成を検出するシステムを構築・検証してきたが、実環境での運用やさらなる精度向上に向けては以下の課題が残されている。

\begin{enumerate}
    \item \textbf{極度の混雑・高密度環境における対応強化}：
    イベント会場やテーマパークのピーク時など、人物同士の接触・重なり（重度のオクルージョン）が定常的に発生する環境においては、YOLOv12 Nano やカルマンフィルタをもってしても一時的な識別IDの消失や追跡途切れが生じる場合がある。密集状態に特化した軽量なヘッド検出や再識別（Re-ID）処理の統合が今後の検討課題である。

    \item \textbf{カメラの設置アングル・画角変化への自動追従}：
    現状のシステムでは、カメラの設置角度や高さを固定パラメータとして前提としているケースがある。設置場所の制約により俯瞰撮影が困難な場合や、画角の歪みが大きい環境においても、キャリブレーションなしで自動的に列線や距離関係を正規化するアルゴリズムの拡張が求められる。

    \item \textbf{動的環境光や天候・照明変動に対する堅牢性}：
    屋外待機列での直射日光、夜間の照明不足、または逆光環境下においては、バウンディングボックスの境界抽出や姿勢ベクトルの信頼度が低下する傾向がある。様々な照明条件下での学習データの追加学習や、エッジ側での動的なコントラスト調整などの前処理層の強化が挙げられる。

    \item \textbf{マルチカメラ連携による広域・長距離待機列のトラッキング}：
    単一のカメラ画像領域（ROI）を超える広大な待機列や折り返しの多い超長距離の列に対しては、複数のエッジ端末またはカメラ間で追跡IDおよびグループ情報をリアルタイムに同期・統合する分散処理基盤の構築が必要となる。
\end{enumerate}
